\chapter{Eigenvalues and Diagonalization}

Every minute, a network monitor splits all incoming requests into three buckets: host A, host B,
and host C. From old logs it builds a transition rule. Of the requests heading to A, 80 percent
stay at A, 10 percent move to B, and 10 percent move to C. From B, 30 percent move to A, 60
percent stay, and 10 percent move to C. From C, 20 percent move to A, 30 percent to B, and 50
percent stay. If the share of requests in the three buckets at minute $k$ is written as the
column $\mathbf{x}_{k}$, that rule takes the form
$\mathbf{x}_{k+1}=M\mathbf{x}_{k}$ with
\[
  M=\begin{pmatrix}0.8&0.3&0.2\\ 0.1&0.6&0.3\\ 0.1&0.1&0.5\end{pmatrix}.
\]

Starting from $\mathbf{x}_{0}=(1,0,0)$, that is, all requests initially go to A, the monitor
computes ten steps and obtains
$\mathbf{x}_{10}=(0.5672,\;0.2661,\;0.1666)$. Thirty steps later the numbers are
$(0.5667,\;0.2667,\;0.1667)$. The three entries still sum to one, and on the second reading the
numbers stop moving at four decimal places.

The question worth asking is not what the result is after forty minutes; repeated multiplication
has already answered that for one starting point. The more useful question is: is there a vector
that $M$ only enlarges without changing its direction? Here that direction is
$(17/30,\;4/15,\;1/6)=(0.5667,\;0.2667,\;0.1667)$, exactly the numbers that appear after forty
steps, and the multiplier is $1$.

Chapter 8 supplied the determinant: $\det M=1/5$, so $M$ is invertible and every system
$M\mathbf{y}=\mathbf{b}$ has a unique solution. That answers whether the process can be reversed,
not which way it is heading, and not how fast. The number that governs the rate here is $1/2$: the
largest eigenvalue after $1$. Since $(1/2)^{40}\approx9.1\times10^{-13}$, the initial error stays
behind the first six digits, and that is why the second reading looks still.

This chapter prepares the tools for reading those two numbers straight from the matrix: how to
compute them, the condition under which the preserved directions are numerous enough to form a
basis, what changes when the basis is changed, and then three uses: matrix powers, linear
differential systems, and Markov chains.

\begin{untukanak}
Take a sheet of rubber marked with a square grid, then pull it from two sides. The squares turn
into diamonds, and almost every line becomes slanted. But one line stays straight and stays
parallel to itself: the line along the pull. Its length changes, its direction does not. Sometimes
there is also another line that does not slant but reverses its direction, like an arrow turned
half a circle.

If we are told only how many times each of those special lines stretches, we already know the
whole pull, without having to draw the grid again. That is a very economical method: three or four
numbers replace thousands of points on the sheet.

The trouble is that not every pull is that tidy. Try twisting the sheet slightly while pulling.
Now no line is really straight any longer, and the economical method fails completely. Some pulls
can be described entirely by a few stretch numbers, others cannot.

The question: from the picture of the pull alone, can we guess whether there are special lines
that do not slant, and how many? In this chapter we will learn why some pulls have such lines and
some do not, why some pulls can be replaced entirely by just a few numbers, and why a computing
engine sometimes guesses wrong on pulls that are almost a twist.
\end{untukanak}

\section{Directions that do not rotate}
\tingkat{basic}

A square matrix acts on vectors: it rotates, stretches, or flattens. Almost every vector comes out
with a direction different from the one it went in with. A small minority do not. A vector that
comes out in the same direction it went in, only with a different length, marks a direction that
the matrix preserves. It is that kind of direction that we name first, before any computing tool
exists.

\kamus{eigenvalue}{nilai eigen}{The multiplier along a direction that is not rotated: the direction
stays the same, its length is multiplied by that number.}

\begin{definisi}[Eigenvalues and eigenvectors]
Let $A$ be an $n\times n$ matrix with real entries. A scalar $\lambda$ is called an
\emph{eigenvalue} of $A$ if there is a nonzero vector $\mathbf{v}\in\mathbb{R}^{n}$ with
$A\mathbf{v}=\lambda\mathbf{v}$. The vector $\mathbf{v}$ is called an \emph{eigenvector} for
$\lambda$.\index{eigenvalue}\index{eigenvector}
\end{definisi}

The condition $\mathbf{v}\neq\mathbf{0}$ is not decoration. Without it, $\mathbf{0}$ satisfies
$A\mathbf{0}=\lambda\mathbf{0}$ for every $\lambda$, so the term loses its meaning. What we are
looking for is precisely a direction, and the zero point has no direction.

The equation $A\mathbf{v}=\lambda\mathbf{v}$ can be rewritten as
$(A-\lambda I)\mathbf{v}=\mathbf{0}$. So $\lambda$ is an eigenvalue exactly when the matrix
$A-\lambda I$ has a nontrivial kernel, that is, exactly when $A-\lambda I$ is singular. A question
about eigenvalues can therefore always be translated into a question about rank deficiency, and
Chapter 8 already supplied the tool for that.

Before computing, one important warning. Not every real matrix has real eigenvalues. The rotation
matrix
\[
  R=\begin{pmatrix}0&-1\\ 1&0\end{pmatrix}
\]
sends $(1,0)$ to $(0,1)$. Every real vector rotates ninety degrees, so no line is preserved;
$R^{2}=-I$ turns half a circle after two steps. The absence of real eigenvalues here is not a
failure of computation but a geometric statement.

To find eigenvalues without guessing, we use the determinant: the eigenvalues are the numbers that
make the matrix lose its inverse once that number is subtracted from the diagonal.

\kamus{characteristic polynomial}{polinomial karakteristik}{The list of numbers that make the
matrix singular, written as a polynomial in one variable.}

\begin{definisi}[Characteristic polynomial]
For an $n\times n$ matrix, the \emph{characteristic polynomial} is defined as
$p_{A}(\lambda)=\det(\lambda I-A)$.\index{characteristic polynomial}
\end{definisi}

We use $\lambda I-A$, not $A-\lambda I$. The two differ by a factor $(-1)^{n}$, so their roots are
the same, but the first choice makes the polynomial monic, and that makes the other coefficients
easier to read, as the next corollary shows.

\begin{teorema}
A scalar $\lambda$ is an eigenvalue of $A$ if and only if $p_{A}(\lambda)=0$.
\end{teorema}

\begin{proof}
Suppose $A\mathbf{v}=\lambda\mathbf{v}$ with $\mathbf{v}\neq\mathbf{0}$. Then
$(A-\lambda I)\mathbf{v}=\mathbf{0}$, so $A-\lambda I$ has a nontrivial kernel and is therefore
singular. By Chapter 8, singularity is equivalent to a zero determinant, so
$\det(A-\lambda I)=0$, which differs by a factor $(-1)^{n}$ from $\det(\lambda I-A)=0$.

Conversely, if $\det(\lambda I-A)=0$ then $\lambda I-A$ is singular. Its kernel is nontrivial, so
there is $\mathbf{v}\neq\mathbf{0}$ with $(\lambda I-A)\mathbf{v}=\mathbf{0}$, that is,
$A\mathbf{v}=\lambda\mathbf{v}$.
\end{proof}

The polynomial $p_{A}$ has degree $n$ and is monic, so it has at most $n$ distinct roots. Over
$\mathbb{C}$ it always factors into $n$ linear factors, so every complex square matrix has exactly
$n$ eigenvalues when counted with multiplicity. Over $\mathbb{R}$ that can fail:
$p_{R}(\lambda)=\lambda^{2}+1$ has no real root, but has a pair of complex roots $\pm i$.

The top two coefficients of the characteristic polynomial carry two quantities we already know.

\begin{akibat}
With eigenvalues counted together with their multiplicities,
$\operatorname{tr}(A)=\sum_{i}\lambda_{i}$ and $\det(A)=\prod_{i}\lambda_{i}$.
\index{trace}
\end{akibat}

\begin{proof}
The coefficient of $\lambda^{n-1}$ in $\det(\lambda I-A)$ is $-\operatorname{tr}(A)$. That term is
produced only by the identity permutation in the Leibniz formula, because every other permutation
loses at least two diagonal entries. The constant term is
$p_{A}(0)=\det(-A)=(-1)^{n}\det(A)$. Comparing those two coefficients with the form
$\prod_{i}(\lambda-\lambda_{i})=\lambda^{n}-(\sum_{i}\lambda_{i})\lambda^{n-1}+\dots+(-1)^{n}
\prod_{i}\lambda_{i}$ gives both identities.
\end{proof}

That corollary gives a cheap check: after computing eigenvalues with great effort, their sum and
their product must match the trace and the determinant. If they do not, there is an arithmetic
mistake.

\begin{bacakeras}
Two voices, taking turns.

First voice: what does a matrix do?\\
Second voice: it rotates, stretches, flattens.\\
First voice: to whom?\\
Second voice: to almost every vector.\\
First voice: and then some do not rotate?\\
Second voice: some do not. One direction. Sometimes two. Never more than the number of columns.\\
First voice: does the length change?\\
Second voice: it changes. The direction stays.\\
First voice: by how much?\\
Second voice: that is the number. That number is the eigenvalue.\\
First voice: and the direction?\\
Second voice: that is the eigenvector. One direction, one number.\\
First voice: what if the direction flips?\\
Second voice: the number is negative. The line stays the same, it only faces backwards.\\
First voice: what if the number is zero?\\
Second voice: that direction disappears, and the matrix loses its inverse.

Together: a direction that does not change stores a number that changes its length.

Repeat once in a normal voice: a direction that does not change stores a number that changes its
length.
\end{bacakeras}

\begin{resep}
For $A=\begin{pmatrix}a&b\\c&d\end{pmatrix}$ whose entries are already known:
\begin{enumerate}
  \item Subtract $\lambda$ from both diagonal entries, leaving $A-\lambda I$.
  \item Compute $\det(A-\lambda I)$, or use $\lambda^{2}-(a+d)\lambda+(ad-bc)$ directly.
  \item Solve that quadratic equation; check that the sum of the roots is $a+d$ and the product is
        $ad-bc$.
  \item For each root $\lambda$, solve $(A-\lambda I)\mathbf{v}=\mathbf{0}$ and take one nonzero
        solution.
\end{enumerate}
Steps 3 and 4 must not be swapped: an eigenvector only has meaning after $\lambda$ is known, and
the column order of $P$ must follow the diagonal order of $D$.
\end{resep}

Now one example by hand. Take
\[
  A=\begin{pmatrix}4&2\\ 1&3\end{pmatrix},\qquad
  p_{A}(\lambda)=\lambda^{2}-7\lambda+10=(\lambda-5)(\lambda-2).
\]
The trace is $7=5+2$ and the determinant is $10=5\cdot2$, which matches. For $\lambda=5$ we solve
$(A-5I)\mathbf{v}=\mathbf{0}$, that is $-v_{1}+2v_{2}=0$ and $v_{1}-2v_{2}=0$; both rows say the
same thing, so $\mathbf{v}=(2,1)$. For $\lambda=2$ we obtain $2v_{1}+2v_{2}=0$, so
$\mathbf{v}=(1,-1)$. Direct check:
\[
  A\begin{pmatrix}2\\ 1\end{pmatrix}=\begin{pmatrix}10\\ 5\end{pmatrix}=5
  \begin{pmatrix}2\\ 1\end{pmatrix},\qquad
  A\begin{pmatrix}1\\ -1\end{pmatrix}=\begin{pmatrix}2\\ -2\end{pmatrix}=2
  \begin{pmatrix}1\\ -1\end{pmatrix}.
\]
Both of those checks can be seen in the plane. Figure~\ref{fig:panah-eigen} draws both eigen
directions together with their images, and one vector that is not an eigenvector as a comparison.

\begin{figure}[htbp]
\centering
\begin{tikzpicture}[x=0.4cm,y=0.4cm,>=latex]
  \draw[->,layergrey!60] (-0.8,0) -- (11.4,0) node[right,font=\scriptsize] {$x$};
  \draw[->,layergrey!60] (0,-2.9) -- (0,6.4) node[above,font=\scriptsize] {$y$};
  \draw[->,thick,layergrey] (0,0) -- (1,0);
  \draw[->,thick,dashed,layergrey] (0,0) -- (4,1)
        node[right,font=\scriptsize] {$(4,1)$};
  \draw[->,very thick,layerteal] (0,0) -- (2,1)
        node[above left,font=\scriptsize] {$(2,1)$};
  \fill[layerteal] (2,1) circle (2pt);
  \draw[->,very thick,accent] (0,0) -- (10,5)
        node[above,font=\scriptsize] {$(10,5)=5\,(2,1)$};
  \draw[->,very thick,layerindigo] (0,0) -- (1,-1)
        node[below left,font=\scriptsize] {$(1,-1)$};
  \fill[layerindigo] (1,-1) circle (2pt);
  \draw[->,very thick,layerindigo] (0,0) -- (2,-2)
        node[below right,font=\scriptsize] {$(2,-2)=2\,(1,-1)$};
\end{tikzpicture}
\caption{Eigen directions of the matrix $A=\begin{pmatrix}4&2\\1&3\end{pmatrix}$. The vector
$(2,1)$ comes out as $(10,5)$, that is, $5$ times as long without rotating; the vector $(1,-1)$
comes out as $(2,-2)$, that is, $2$ times as long. The vector $(1,0)$ comes out as $(4,1)$ and
rotates, so it is not an eigenvector.}
\label{fig:panah-eigen}
\end{figure}

An eigenvector is determined only up to a multiplier: $(4,2)$ is also an eigenvector for $5$. What
is unique is the direction, not the length. We will use that freedom later to tidy up the matrix.

For a $2\times2$ matrix the roots can be written down directly. Since
$p_{A}(\lambda)=\lambda^{2}-\operatorname{tr}(A)\lambda+\det(A)$, the quadratic formula gives
\[
  \lambda_{1,2}=\frac{\operatorname{tr}(A)\pm\sqrt{\operatorname{tr}(A)^{2}-4\det(A)}}{2},
\]
so the sign of $\Delta=\operatorname{tr}(A)^{2}-4\det(A)$ determines the whole spectrum. For $A$
above, $\Delta=49-40=9$, and the roots are $5$ and $2$. If $\Delta=0$ the two roots coincide, and
if $\Delta<0$ there is no real root: the rotation matrix above gives $\Delta=0-4=-4$.

\begin{figure}[htbp]
\centering
\begin{tikzpicture}[x=0.85cm,y=0.85cm,>=latex,font=\scriptsize]
  \draw[->,layergrey!60] (-1.2,0.35) -- (6.6,0.35);
  \foreach \x in {0,1,2,3,4,5,6} \draw[layergrey!50] (\x,0.26) -- (\x,0.44);
  \fill[layerteal] (2,0.35) circle (2.2pt);
  \fill[layerteal] (5,0.35) circle (2.2pt);
  \node[above=2pt] at (2,0.35) {$2$};
  \node[above=2pt] at (5,0.35) {$5$};
  \node[left,align=right,text width=3.0cm] at (-1.3,0.35) {$\Delta=9$: two real roots};

  \draw[->,layergrey!60] (-1.2,-1.25) -- (6.6,-1.25);
  \foreach \x in {0,1,2,3,4,5,6} \draw[layergrey!50] (\x,-1.34) -- (\x,-1.16);
  \fill[layerorange] (1,-1.25) circle (2.2pt);
  \node[above=2pt] at (1,-1.2) {$1$};
  \node[left,align=right,text width=3.0cm] at (-1.3,-1.25) {$\Delta=0$: one double root};

  \draw[->,layergrey!60] (-1.2,-2.85) -- (6.6,-2.85);
  \foreach \x in {0,1,2,3,4,5,6} \draw[layergrey!50] (\x,-2.94) -- (\x,-2.76);
  \draw[layergrey!80] (0,-3.55) -- (0,-2.15);
  \draw[accent,fill=white,thick] (0,-2.3) circle (2.4pt);
  \draw[accent,fill=white,thick] (0,-3.4) circle (2.4pt);
  \node[left=2pt,accent] at (0,-2.3) {$i$};
  \node[left=2pt,accent] at (0,-3.4) {$-i$};
  \node[left,align=right,text width=3.0cm] at (-1.3,-2.85) {$\Delta=-4$: no real roots};
\end{tikzpicture}
\caption{The spectrum of a $2\times2$ matrix on the number line, separated by the sign of
$\Delta=\operatorname{tr}(A)^{2}-4\det(A)$. Top row: $A$ with $\Delta=49-40=9$ gives the roots
$2$ and $5$. Middle row: a shear matrix with $\Delta=0$ gives the root $1$ counted twice. Bottom
row: the rotation matrix with $\Delta=0-4=-4$ does not meet the real line, and both of its roots,
$\pm i$, lie off that line.}
\label{fig:spektrum}
\end{figure}

Figure~\ref{fig:spektrum} arranges those three possibilities on a single number line: two points,
one point counted twice, or no point at all. The last is not a shortcoming of the picture but a
statement that the real line contains no preserved direction.

\begin{tangan}
Materials: one rubber band, two pencils, a sheet of paper, a ruler.

Place the two pencils parallel, then stretch the rubber band around both so that a rectangle
forms. On the paper, draw three arrows from the lower left corner of that rectangle, each at an
angle of 30, 45, and 60 degrees to the bottom side. Label all three.

Slide the right pencil twice as far, so that the rubber stretches. Compare each arrow with its
original length. The arrow parallel to the pull side keeps its direction; the other two slant, and
the one that slants most is the one farthest from the pull side.

Now rotate the whole frame ninety degrees and repeat. The surviving direction rotates along with
the rubber, because what is special is the direction relative to the rubber, not relative to the
table.

What should be visible: one direction survives, two directions rotate, and the stretch factor
along the surviving direction is the multiplier we call the eigenvalue.
\end{tangan}

\section{Eigen spaces and two multiplicities}
\tingkat{intermediate}

One eigenvalue usually comes with more than one eigenvector. The set of all eigenvectors for one
eigenvalue, plus the zero vector, is not a random collection: it is a subspace. Its dimension will
later decide whether the matrix can be simplified.

\begin{definisi}[Eigen space]
For an eigenvalue $\lambda$ of $A$, the \emph{eigen space} is
$E_{\lambda}=\operatorname{null}(A-\lambda I)=\{\mathbf{v}:A\mathbf{v}=\lambda\mathbf{v}\}$.
\index{eigen space}
\end{definisi}

\kamus{eigen space}{ruang eigen}{All directions that the matrix does not rotate, gathered into a
single subspace.}

An eigen space is always a subspace because it is the kernel of a matrix (Chapter 7). Its
dimension is at least one exactly when $\lambda$ really is an eigenvalue.

Eigen spaces for different eigenvalues intersect only at the zero point.

\begin{proposisi}
If $\lambda\neq\mu$, then $E_{\lambda}\cap E_{\mu}=\{\mathbf{0}\}$.
\end{proposisi}

\begin{proof}
Take $\mathbf{v}$ in that intersection. Then $A\mathbf{v}=\lambda\mathbf{v}$ and
$A\mathbf{v}=\mu\mathbf{v}$, so $(\lambda-\mu)\mathbf{v}=\mathbf{0}$. Since
$\lambda-\mu\neq0$, we must have $\mathbf{v}=\mathbf{0}$.
\end{proof}

That property extends to any finite number of eigenvectors, by induction. The result is the lemma
that forms the backbone of this chapter.

\begin{lemma}
Eigenvectors corresponding to distinct eigenvalues are linearly independent of one another.
\end{lemma}

\begin{proof}
We prove this by induction on the number of vectors. For one vector the statement is true because
an eigenvector is nonzero. Suppose the statement holds for $m-1$ vectors, and let
\[
  c_{1}\mathbf{v}_{1}+\dots+c_{m}\mathbf{v}_{m}=\mathbf{0}
\]
with $\lambda_{1},\dots,\lambda_{m}$ distinct. Multiply that equation by $A$, and use
$A\mathbf{v}_{i}=\lambda_{i}\mathbf{v}_{i}$:
\[
  c_{1}\lambda_{1}\mathbf{v}_{1}+\dots+c_{m}\lambda_{m}\mathbf{v}_{m}=\mathbf{0}.
\]
Subtract $\lambda_{m}$ times the first equation from the second:
\[
  c_{1}(\lambda_{1}-\lambda_{m})\mathbf{v}_{1}+\dots
  +c_{m-1}(\lambda_{m-1}-\lambda_{m})\mathbf{v}_{m-1}=\mathbf{0}.
\]
By the induction hypothesis, $\mathbf{v}_{1},\dots,\mathbf{v}_{m-1}$ are independent, so
$c_{i}(\lambda_{i}-\lambda_{m})=0$ for $i<m$. Since $\lambda_{i}\neq\lambda_{m}$, all $c_{i}=0$
for $i<m$. The original equation is left with $c_{m}\mathbf{v}_{m}=\mathbf{0}$, so $c_{m}=0$.
\end{proof}

That lemma gives a tidy upper bound: an $n\times n$ matrix cannot have more than $n$ independent
eigenvectors. When is that bound attained?

To answer, we need two ways of counting how often one eigenvalue occurs: from the polynomial and
from the dimension of the eigen space. The two often differ, and it is the difference that decides.

\kamus{multiplicity}{multiplisitas}{How many times an eigenvalue occurs, counted in two ways that
are not always equal.}

\begin{definisi}[Algebraic and geometric multiplicity]
Let $\lambda$ be an eigenvalue of $A$. The \emph{algebraic multiplicity} of $\lambda$ is the power
of the factor $(\lambda'-\lambda)$ in $p_{A}(\lambda')$. The \emph{geometric multiplicity} of
$\lambda$ is $\dim E_{\lambda}=\operatorname{nullity}(A-\lambda I)$.\index{multiplicity}
\end{definisi}

Geometric multiplicity is easy to compute: subtract $\lambda$ from the diagonal, then compute the
nullity of the resulting matrix. Algebraic multiplicity demands factoring a polynomial. The
relation between the two always runs in one direction.

\begin{lemma}
For every eigenvalue $\lambda$ we have
$1\le\dim E_{\lambda}\le$ the algebraic multiplicity of $\lambda$.
\end{lemma}

\begin{proof}
The left inequality is true because $\lambda$ is an eigenvalue, so $E_{\lambda}$ contains a
nonzero vector.

For the right inequality, write $k=\dim E_{\lambda}$ and choose a basis
$\mathbf{u}_{1},\dots,\mathbf{u}_{k}$ of $E_{\lambda}$. Extend it to a basis
$\mathbf{u}_{1},\dots,\mathbf{u}_{n}$ of $\mathbb{R}^{n}$, which is possible by Chapter 7. In that
basis, column $j$ of the matrix $A$ for $j\le k$ is $A\mathbf{u}_{j}=\lambda\mathbf{u}_{j}$, that
is, $\lambda$ in the diagonal position and zero everywhere else. So the matrix $A$ in that basis
has the form
\[
  \begin{pmatrix}\lambda I_{k}&C\\ 0&D\end{pmatrix}
\]
with $C$ of size $k\times(n-k)$ and $D$ of size $(n-k)\times(n-k)$. The determinant of
$\lambda' I-A$ in that basis, either through cofactor expansion along the first $k$ columns or
through the determinant formula for a block triangular matrix, equals
$(\lambda'-\lambda)^{k}\det(\lambda' I-D)$. So $(\lambda'-\lambda)^{k}$ divides $p_{A}$, hence the
algebraic multiplicity of $\lambda$ is at least $k$.
\end{proof}

Now look at both possibilities in a single matrix. Take
\[
  A=\begin{pmatrix}2&1&1\\ 1&2&1\\ 1&1&2\end{pmatrix},\qquad
  p_{A}(\lambda)=(\lambda-1)^{2}(\lambda-4).
\]
The eigenvalues are $1$ and $4$, with algebraic multiplicities $2$ and $1$. The trace is
$6=4+1+1$ and the determinant is $4=4\cdot1\cdot1$, which matches. For $\lambda=4$, the equation
$(A-4I)\mathbf{v}=\mathbf{0}$ gives $v_{1}=v_{2}=v_{3}$, so $E_{4}$ is spanned by $(1,1,1)$. For
$\lambda=1$, the matrix $A-I$ is the matrix all of whose entries are $1$, so its kernel is the set
of vectors whose entries sum to zero. That set is two-dimensional, with basis $(1,-1,0)$ and
$(1,0,-1)$. So the geometric multiplicity of $1$ is $2$, equal to its algebraic multiplicity.

A repeated eigenvalue is not automatically a problem; what becomes a problem is an eigen space
that is too small, and that happens with the following matrix.

\section{When a matrix can be diagonalized}
\tingkat{intermediate}

If we may choose the basis, we want the matrix as simple as possible: diagonal, so that every
basis vector is merely scaled and not mixed with the others. Is there always such a basis?

\begin{definisi}[Diagonalizable]
An $n\times n$ matrix is called \emph{diagonalizable} if there is an invertible matrix $P$ and a
diagonal matrix $D$ with $P^{-1}AP=D$, or equivalently $A=PDP^{-1}$.
\end{definisi}

\kamus{diagonalization}{diagonalisasi}{Changing the measuring instrument until the whole work of
the matrix is a few multiplications of numbers.}

The columns of $P$ will contain eigenvectors, and $D$ contains the eigenvalues on its diagonal. A
question about diagonalization is therefore really a question about having enough eigenvectors.

\begin{teorema}
An $n\times n$ matrix is diagonalizable if and only if it has $n$ linearly independent
eigenvectors.\index{diagonalization}
\end{teorema}

\begin{proof}
Suppose $P^{-1}AP=D$ with $D=\operatorname{diag}(d_{1},\dots,d_{n})$. Multiply from the left by
$P$: $AP=PD$. Column $j$ of both sides gives
$A\mathbf{p}_{j}=d_{j}\mathbf{p}_{j}$, with $\mathbf{p}_{j}$ the $j$-th column of $P$. Since $P$ is
invertible, each of its columns is nonzero, so $\mathbf{p}_{j}$ is an eigenvector. The columns of
$P$ are also linearly independent because $P$ is invertible. So $A$ has $n$ independent
eigenvectors.

Conversely, suppose $\mathbf{p}_{1},\dots,\mathbf{p}_{n}$ are independent eigenvectors with
$A\mathbf{p}_{j}=\lambda_{j}\mathbf{p}_{j}$. Assemble $P$ from those vectors as columns and
$D=\operatorname{diag}(\lambda_{1},\dots,\lambda_{n})$. Then $AP=PD$, and since $P$ is invertible
we obtain $P^{-1}AP=D$.
\end{proof}

The corollary most often used in hand computation: enough distinct eigenvalues is already
sufficient.

\begin{akibat}
If an $n\times n$ matrix has $n$ distinct eigenvalues, it is diagonalizable.
\end{akibat}

\begin{proof}
Take one eigenvector for each eigenvalue. By the earlier lemma, $n$ eigenvectors with $n$ distinct
eigenvalues are independent of one another. The theorem above then gives a diagonalization.
\end{proof}

The converse does not hold: the $3\times3$ example above had only two distinct eigenvalues, yet is
still diagonalizable because the eigen space for $\lambda=1$ is two-dimensional. The exact
condition is not the number of eigenvalues but the total dimension of the eigen spaces.

\begin{teorema}[Multiplicity condition]
An $n\times n$ matrix is diagonalizable if and only if for each of its eigenvalues the algebraic
multiplicity equals the geometric multiplicity. That condition is equivalent to
$\sum_{\lambda}\dim E_{\lambda}=n$.
\end{teorema}

\begin{proof}
First, suppose the total dimension of all eigen spaces equals $n$. For each $\lambda$ choose a
basis of $E_{\lambda}$, then combine all those bases into a list $\mathbf{v}_{1},\dots,
\mathbf{v}_{n}$. That list is linearly independent: if some linear combination equals zero, group
the terms according to their eigenvalues to get $\mathbf{u}_{1}+\dots+\mathbf{u}_{m}=\mathbf{0}$
with $\mathbf{u}_{k}\in E_{\lambda_{k}}$. By the lemma on distinct eigenvectors, every
$\mathbf{u}_{k}=\mathbf{0}$; since each $\mathbf{u}_{k}$ is a combination of the basis of
$E_{\lambda_{k}}$, all its coefficients are zero. That list also contains $n$ vectors in a space of
dimension $n$, so it is a basis of eigenvectors, and the previous theorem closes the argument.

Second, suppose $A$ is diagonalizable, so there is a basis of eigenvectors
$\mathbf{p}_{1},\dots,\mathbf{p}_{n}$. In that basis the matrix $A$ is diagonal, so
$p_{A}(\lambda')=\prod_{j}(\lambda'-\lambda_{j})$ and the algebraic multiplicity of each
eigenvalue equals the number of $j$ with $\lambda_{j}=\lambda$. That number of $j$ does not exceed
$\dim E_{\lambda}$, because those vectors lie inside $E_{\lambda}$ and are independent. So the
algebraic multiplicity does not exceed the geometric one. The earlier lemma gives the opposite
direction, so the two are equal. Finally, $\sum_{\lambda}\dim E_{\lambda}=n$ because the basis of
eigenvectors fills the whole space.
\end{proof}

Now an example that fails. Take
\[
  S=\begin{pmatrix}1&1\\ 0&1\end{pmatrix},\qquad p_{S}(\lambda)=(\lambda-1)^{2}.
\]
The only eigenvalue is $1$, with algebraic multiplicity $2$. But
$S-I=\begin{pmatrix}0&1\\ 0&0\end{pmatrix}$ has rank one, so $\dim E_{1}=2-1=1$. Its geometric
multiplicity is one short, and by the theorem above $S$ is not diagonalizable.

Geometrically $S$ is a shear: it preserves the direction $(1,0)$ and slides $(0,1)$ to $(1,1)$. No
basis makes this map a mere scaling, because a scaling never slides something that was previously
parallel. That failure is a property of the map, not a choice of basis.

\begin{untukanak}
Have you ever tried to tidy something that turned out not to be tidyable? Take a stack of cards,
then slide the top part slightly sideways. The picture on each card stays the same, but every
point shifts by its distance from the base. There is no way to rotate that stack so that the shift
becomes a mere lengthening or shortening of something. A shift remains a shift.

There are other objects that are more accommodating. A rubber sheet pulled straight can be
described only by how many times each direction stretches. If we know that number for a few
important directions, we know the whole pull, including for points we never measured.

The question: how do we tell the two apart from the picture alone, without trying every
possibility? In this chapter we will learn why a shear always stores one direction that cannot be
tidied up, why some pulls have enough special directions, and why that difference decides whether
the tenth power from a computing engine can still be trusted.
\end{untukanak}

A shear is easy to recognize because of the triangle. The same failure appears in matrices that do
not look triangular. Take
\[
  B=\begin{pmatrix}-1&4&0\\ -1&3&0\\ 1&-2&2\end{pmatrix},\qquad
  p_{B}(\lambda)=(\lambda-1)^{2}(\lambda-2).
\]
For $\lambda=1$, the matrix $B-I$ has rank two, so $\dim E_{1}=3-2=1$, while its algebraic
multiplicity is $2$. The only eigenvector for $1$ is $(2,1,0)$. For $\lambda=2$ its eigen space is
spanned by $(0,0,1)$ and is one-dimensional, matching its algebraic multiplicity. The total
dimension of all eigen spaces is only $2<3$, so $B$ is not diagonalizable. A matrix of this kind is
called \emph{defective}: it is short of eigenvectors because its eigen spaces are narrow.
\index{defective}

\kamus{defective}{defective}{A matrix whose eigenvalues repeat but whose eigen spaces are too
narrow to form a basis.}

One note before we go on. Whether a matrix is diagonalizable depends on the field of scalars: $R$
is not diagonalizable over $\mathbb{R}$, but over $\mathbb{C}$ it has two distinct eigenvalues,
$\pm i$, so it is diagonalizable. A statement without a field qualifier is therefore incomplete.
For the rest of this chapter the field is $\mathbb{R}$ unless stated otherwise.

\section{Similar matrices and invariants}
\tingkat{advanced}

Diagonalization is a change of basis that produces the simplest possible form. But not every
matrix can reach that form, and the next question is natural: what stays the same across all
matrices connected by a change of basis?\index{invariant}

\begin{definisi}[Similar matrices]
Two $n\times n$ matrices $A$ and $B$ are called \emph{similar} if there is an invertible matrix
$P$ with $B=P^{-1}AP$.\index{similar}
\end{definisi}

\kamus{similar}{serupa}{Two reports about the same map from two different measuring spaces.}

Similarity is an equivalence relation: every matrix is similar to itself, the relation is
symmetric because $(P^{-1})^{-1}BP^{-1}=A$, and transitive because
$(PQ)^{-1}A(PQ)=Q^{-1}(P^{-1}AP)Q$. The classes divide up neatly, and diagonalization means asking
whether a class contains a diagonal matrix.

\begin{proposisi}
If $A$ and $B$ are similar, then $p_{A}=p_{B}$. Consequently $A$ and $B$ have the same eigenvalues
together with their algebraic multiplicities, and also the same trace, determinant, rank, and
nullity.
\end{proposisi}

\begin{proof}
Suppose $B=P^{-1}AP$. Then
\[
  \lambda I-B=\lambda I-P^{-1}AP=P^{-1}(\lambda I-A)P,
\]
so
$\det(\lambda I-B)=\det(P^{-1})\det(\lambda I-A)\det(P)=\det(\lambda I-A)$ because
$\det(P^{-1})\det(P)=1$. So the characteristic polynomials are equal, and all data derived from
their coefficients are equal as well. Equality of rank and nullity follows from $B-\lambda I=
P^{-1}(A-\lambda I)P$: multiplication by an invertible matrix does not change the rank, and it is
that rank which determines the nullity through rank-nullity.
\end{proof}

The converse does not hold, and the counterexample is simple: $I_{2}$ and the shear $S$ have the
same characteristic polynomial, $(\lambda-1)^{2}$, but are not similar. If $P^{-1}I_{2}P=S$, the
left side equals $I_{2}$ for every $P$. So the characteristic polynomial is not a complete
invariant; it only filters.

What changes and what does not can be separated neatly. The eigenvalues do not change. The
eigenvectors do change, and they change in a way that is easy to remember.

\begin{proposisi}
If $B=P^{-1}AP$ and $A\mathbf{v}=\lambda\mathbf{v}$, then
$B(P^{-1}\mathbf{v})=\lambda(P^{-1}\mathbf{v})$.
\end{proposisi}

\begin{proof}
Compute $B(P^{-1}\mathbf{v})=P^{-1}APP^{-1}\mathbf{v}=P^{-1}A\mathbf{v}
=P^{-1}(\lambda\mathbf{v})=\lambda(P^{-1}\mathbf{v})$.
\end{proof}

So eigenvectors move along with their coordinates, while the multipliers stay put. That makes
sense: a preserved direction is a property of the map, while its list of coordinates depends on
the basis. That is why eigenvalues are used as the fingerprint of a matrix, and eigenvectors are
not.

On the other hand, the freedom to choose $P$ explains why two diagonalizations can look different
although both are correct: the columns of $P$ may be multiplied by a nonzero scalar or permuted as
long as $D$ is permuted accordingly. What must be the same is the set of eigenvalues together with
their multiplicities.

\section{A worked example}
\tingkat{intermediate}

Every number in this section is computed by hand first, then checked with Python in the lab
section. We start from the $3\times3$ matrix that already appeared in the discussion of eigen
spaces,
\[
  A=\begin{pmatrix}2&1&1\\ 1&2&1\\ 1&1&2\end{pmatrix}.
\]
Its characteristic polynomial is computed from $\lambda I-A$, whose diagonal entries are all
$\lambda-2$ and whose other entries are all $-1$. Writing $t=\lambda-2$ and expanding along the
first row,
\[
  p_{A}(\lambda)=t\,(t^{2}-1)-(-1)\,(-t-1)+(-1)\,(1+t)=t^{3}-3t-2 .
\]
That form factors because $t=-1$ is a root: $t^{3}-3t-2=(t+1)^{2}(t-2)$. Back to $\lambda$,
\[
  p_{A}(\lambda)=(\lambda-1)^{2}(\lambda-4)=\lambda^{3}-6\lambda^{2}+9\lambda-4 .
\]
The trace is $6$ and the determinant is $4$, matching $4+1+1$ and $4\cdot1\cdot1$.

Eigenvectors for $\lambda=4$: the matrix $A-4I$ has diagonal entries $-2$ and other entries $1$.
Its rows say $-2v_{1}+v_{2}+v_{3}=0$ and so on; all three equations hold exactly when
$v_{1}=v_{2}=v_{3}$. So $E_{4}$ is spanned by $(1,1,1)$.

Eigenvectors for $\lambda=1$: the matrix $A-I$ contains only $1$s, so its kernel is the set of
vectors with $v_{1}+v_{2}+v_{3}=0$. A basis is two vectors, for instance $(1,-1,0)$ and
$(1,0,-1)$. Hence
\[
  P=\begin{pmatrix}1&1&1\\ 1&-1&0\\ 1&0&-1\end{pmatrix},\qquad
  P^{-1}=\frac{1}{3}\begin{pmatrix}1&1&1\\ 1&-2&1\\ 1&1&-2\end{pmatrix},\qquad
  D=\operatorname{diag}(4,1,1).
\]
The determinant of $P$ is $3$, and that is the denominator in $P^{-1}$. We check one entry by
hand: the entry in the first row and first column of $PDP^{-1}$ is
$\tfrac{1}{3}(4\cdot1\cdot1+1\cdot1\cdot1+1\cdot1\cdot1)=2$, matching the entry of $A$. A full
check with SymPy gives $PDP^{-1}-A=0$ exactly.

Now the tenth power.\index{matrix power} Multiplying nine times gives no insight at all;
diagonalization provides it. Since $A=I+J$ with $J$ the matrix all of whose entries are $1$, and
$J^{2}=3J$, every power of $J$ returns as a multiple of $J$:
\[
  J^{2}=3J,\quad J^{3}=9J,\quad\dots,\quad J^{k}=3^{k-1}J .
\]
Then
\[
  A^{k}=(I+J)^{k}=I+\Bigl(\sum_{j=1}^{k}\binom{k}{j}3^{j-1}\Bigr)J
  =I+\frac{4^{k}-1}{3}\,J,
\]
with the last step using the geometric series formula
$\sum_{j=1}^{k}\binom{k}{j}3^{j-1}=\bigl((1+3)^{k}-1\bigr)/3$. For $k=10$ we have
$4^{10}=1048576$ and $(1048576-1)/3=349525$, so that
\[
  A^{10}=I+349525\,J=
  \begin{pmatrix}349526&349525&349525\\ 349525&349526&349525\\ 349525&349525&349526\end{pmatrix}.
\]
Every entry can be checked: $1+349525$ on the diagonal, $349525$ off it. Direct multiplication
with Python gives the same matrix, with largest difference $1.05\times10^{-9}$ because of
floating-point arithmetic.

The form of that formula explains what is not visible in the final result. The term $4^{k}$
dominates the direction $(1,1,1)$; the constant term $1$ survives in the direction whose entries
sum to zero. The eigenvalues $4$ and $1$ are two growth rates, and a matrix power is merely a
mixture of the two. If the matrix is not diagonalizable, that mixture changes shape: a term
$k\lambda^{k}$ appears or, for $\lambda=1$, a pure factor $k$.

The defective example above shows that. For the $B$ that failed to be diagonalizable,
\[
  B^{k}=
  \begin{pmatrix}1-2k&4k&0\\ -k&1+2k&0\\ 2^{k}-1&-2(2^{k}-1)&2^{k}\end{pmatrix},
\]
which can be proved by induction and holds for $k=1$. At $k=10$ that formula gives
\[
  B^{10}=\begin{pmatrix}-19&40&0\\ -10&21&0\\ 1023&-2046&1024\end{pmatrix},
\]
the same as repeated multiplication. The top-left entry, $-19=1-2\cdot10$, grows linearly, not
exponentially: the eigenvalue $1$ controls that direction, and the shortage of eigenvectors adds
one factor of $k$.

Now the first application: a process that settles. For the Markov chain matrix $M$ from the
opening,
\index{Markov chain}
\[
  p_{M}(\lambda)=\lambda^{3}-\tfrac{19}{10}\lambda^{2}+\tfrac{11}{10}\lambda-\tfrac15
  =(\lambda-1)\bigl(\lambda-\tfrac12\bigr)\bigl(\lambda-\tfrac25\bigr).
\]
The trace is $19/10$ and the determinant is $1/5$, matching $1+\tfrac12+\tfrac25$ and
$1\cdot\tfrac12\cdot\tfrac25$. The eigenvector for $\lambda=1$ is found from $M\mathbf{v}=
\mathbf{v}$, that is $(M-I)\mathbf{v}=\mathbf{0}$, together with the condition that the entries sum
to one. The solution gives $\boldsymbol{\pi}=(17/30,\;4/15,\;1/6)$. One coordinate checked by hand:
\[
  \tfrac{8}{10}\cdot\tfrac{17}{30}+\tfrac{3}{10}\cdot\tfrac{8}{30}
  +\tfrac{2}{10}\cdot\tfrac{5}{30}
  =\tfrac{13.6+2.4+1.0}{30}=\tfrac{17}{30}.
\]
So $\boldsymbol{\pi}$ really is a fixed point.

Since $M$ has three distinct eigenvalues, it is diagonalizable. Every starting point therefore
decomposes as
\[
  \mathbf{x}_{k}=\boldsymbol{\pi}+c_{2}\bigl(\tfrac12\bigr)^{k}\mathbf{v}_{2}
  +c_{3}\bigl(\tfrac25\bigr)^{k}\mathbf{v}_{3},
\]
with $c_{2},c_{3}$ depending on $\mathbf{x}_{0}$. The second and third terms shrink, and the rate
is governed by $1/2$. For $\mathbf{x}_{0}=(1,0,0)$ we obtain $c_{2}\mathbf{v}_{2}$ of order
$10^{-1}$, so the error after ten steps is of order
$10^{-1}\cdot(1/2)^{10}\approx10^{-4}$, and after forty steps of order
$(1/2)^{40}\approx9.1\times10^{-13}$. The computed result agrees: $\mathbf{x}_{10}$ differs from
$\boldsymbol{\pi}$ in the fourth digit, while $\mathbf{x}_{40}$ already agrees to six digits.

Second application: differential systems. If $\mathbf{x}(t)$ satisfies
$\mathbf{x}'=A\mathbf{x}$ with $A$ above, the solution is $\mathbf{x}(t)=e^{At}\mathbf{x}(0)$.
Since $A=I+J$ and $J^{2}=3J$,
\[
  e^{At}=e^{t}e^{Jt}=e^{t}\Bigl(I+\frac{e^{3t}-1}{3}J\Bigr)
  =e^{t}I+\frac{e^{4t}-e^{t}}{3}J .
\]
At $t=0$ that formula gives $I$. For $t=0.5$ its off-diagonal entries are
$(e^{2}-e^{0.5})/3=1.9134$, and Python gives the same number to four digits. The component along
$(1,1,1)$ grows like $e^{4t}$; the component in the plane where the entries sum to zero grows like
$e^{t}$, so the solution quickly becomes parallel to $(1,1,1)$. A positive real eigenvalue means
growth; a purely imaginary eigenvalue such as $\pm i$ means oscillation without growth.

Third application: power iteration. Multiplying repeatedly and then normalizing the result,
$\mathbf{x}_{k+1}=A\mathbf{x}_{k}/\lVert A\mathbf{x}_{k}\rVert$, makes the direction turn toward
the dominant eigenvector, because $\lambda_{1}^{k}$ defeats $\lambda_{2}^{k}$ when
$|\lambda_{1}|>|\lambda_{2}|$. For $A$ above starting from $(1,0)$, the angle runs $0^\circ$,
$14.04^\circ$, $21.25^\circ$, $24.39^\circ$, $25.69^\circ$, $26.21^\circ$, and heads toward
$\arctan(1/2)=26.57^\circ$; the Rayleigh quotient $\mathbf{x}^{\top}A\mathbf{x}/\mathbf{x}^{\top}
\mathbf{x}$ runs $4.00$, $4.65$, $4.88$, $4.96$, $4.98$, $4.99$, heading toward $5$.
Figure~\ref{fig:iterasi-pangkat} shows that turning.

\begin{figure}[htbp]
\centering
\begin{tikzpicture}[scale=3.2,>=latex]
  \draw[layergrey!35] (0,0) circle (1);
  \draw[->,layergrey!60] (-1.25,0) -- (1.55,0) node[right,font=\scriptsize] {$x$};
  \draw[->,layergrey!60] (0,-0.35) -- (0,1.2) node[above,font=\scriptsize] {$y$};
  \draw[->,semithick,layergrey] (0,0) -- (0:1);
  \draw[->,thick,layerindigo!45] (0,0) -- (14.04:1);
  \draw[->,thick,layerindigo!70] (0,0) -- (21.25:1);
  \draw[->,very thick,layerindigo] (0,0) -- (24.39:1);
  \foreach \a in {0,14.04,21.25,24.39} \fill[layerindigo!70] (\a:1) circle (1.3pt);
  \draw[->,very thick,dashed,accent] (0,0) -- (26.565:1.42);
  \node[right,font=\scriptsize,accent,align=left] at (26.565:1.42)
        {$(2,1)$\\[1pt]$26.57^\circ$};
\end{tikzpicture}
\caption{Power iteration on $A=\begin{pmatrix}4&2\\1&3\end{pmatrix}$ starting from $(1,0)$. The
solid arrows are the direction after $k=0,1,2,3$ steps, with angles $0^\circ$, $14.04^\circ$,
$21.25^\circ$, and $24.39^\circ$. The dashed arrow is the dominant eigen direction $(2,1)$ at
$26.57^\circ$, which is approached but never passed; the last two arrows coincide because the
angles are getting tighter, and that is precisely the point of the picture.}
\label{fig:iterasi-pangkat}
\end{figure}

\kertas{For $J$ the matrix all of whose entries are $1$, take $B=J-I$. Determine the eigenvalues of
$B$ first, then predict $B^{10}$ using the pattern $I+\frac{\mu^{10}-(-1)^{10}}{3}J$ with $\mu$ the
nonzero eigenvalue of $B$. Check that prediction with Python, then explain why $B$ is
diagonalizable although it has a repeated eigenvalue. As an additional exercise, compute $e^{Bt}$
and determine the limit $t\to\infty$ of $e^{-2t}e^{Bt}$.}

\section{Computational lab}
\tingkat{advanced}

\begin{lab}
\begin{verbatim}
import numpy as np
from scipy.linalg import expm

np.set_printoptions(precision=4, suppress=True)

A = np.array([[2., 1., 1.],
              [1., 2., 1.],
              [1., 1., 2.]])
lam, V = np.linalg.eig(A)
print("eigenvalues A :", np.round(np.sort(lam.real), 10))
for i in range(3):
    sisa = A @ V[:, i] - lam[i] * V[:, i]
    print("  residual Av-lv :", round(float(np.linalg.norm(sisa)), 16))
print("A^10 via diagonalization :")
print(np.round((V @ np.diag(lam ** 10) @ np.linalg.inv(V)).real).astype(int))
print("A^10 via ordinary power :")
print(np.linalg.matrix_power(A, 10).astype(int))
print("max difference :", np.abs((V @ np.diag(lam ** 10)
                              @ np.linalg.inv(V)).real
                             - np.linalg.matrix_power(A, 10)).max())
print("cond V       :", round(float(np.linalg.cond(V)), 4))

A2 = np.array([[4., 2.],
               [1., 3.]])
x = np.array([1., 0.])
print("power iteration on A2:")
for k in range(6):
    sudut = np.degrees(np.arctan2(x[1], x[0]))
    rq = float(x @ A2 @ x / (x @ x))
    print("  k=%d angle=%7.2f rayleigh=%.4f" % (k, sudut, rq))
    x = A2 @ x / np.linalg.norm(A2 @ x)
print("limit angle  :", round(float(np.degrees(np.arctan2(1., 2.))), 2))

B = np.array([[-1., 4., 0.],
              [-1., 3., 0.],
              [1., -2., 2.]])
print("eigen B      :", np.round(np.sort(np.linalg.eigvals(B).real), 10))
print("rank(B - I)  :", np.linalg.matrix_rank(B - np.eye(3)))
print("B^10         :")
print(np.linalg.matrix_power(B, 10).astype(int))

M = np.array([[0.8, 0.3, 0.2],
              [0.1, 0.6, 0.3],
              [0.1, 0.1, 0.5]])
print("column sums  :", M.sum(axis=0))
print("eigen M      :", np.round(np.sort(np.linalg.eigvals(M).real), 10))
x = np.array([1., 0., 0.])
for _ in range(10):
    x = M @ x
print("x after 10   :", np.round(x, 6))
for _ in range(30):
    x = M @ x
print("x after 40   :", np.round(x, 6))

C = np.array([[1., 1000.], [0., 2.]])
Cp = C.copy()
Cp[1, 0] = 1e-6
print("eigen C      :", np.round(np.sort(np.linalg.eigvals(C).real), 6))
print("eigen C+1e-6 :", np.round(np.sort(np.linalg.eigvals(Cp).real), 6))
print("cond V(C)    :", round(float(np.linalg.cond(np.linalg.eig(C)[1])), 1))
print("expm(0.5A)   :", round(float(expm(0.5 * A)[0, 1]), 4))
print("formula      :", round(float((np.exp(2) - np.exp(0.5)) / 3), 4))
\end{verbatim}
Actual output:
\begin{verbatim}
eigenvalues A : [1. 1. 4.]
  residual Av-lv : 8e-16
  residual Av-lv : 1.1e-15
  residual Av-lv : 3e-16
A^10 via diagonalization :
[[349526 349525 349525]
 [349525 349526 349525]
 [349525 349525 349526]]
A^10 via ordinary power :
[[349526 349525 349525]
 [349525 349526 349525]
 [349525 349525 349526]]
max difference : 1.0477378964424133e-09
cond V       : 1.5227
power iteration on A2:
  k=0 angle=   0.00 rayleigh=4.0000
  k=1 angle=  14.04 rayleigh=4.6471
  k=2 angle=  21.25 rayleigh=4.8820
  k=3 angle=  24.39 rayleigh=4.9578
  k=4 angle=  25.69 rayleigh=4.9840
  k=5 angle=  26.21 rayleigh=4.9938
limit angle  : 26.57
eigen B      : [1. 1. 2.]
rank(B - I)  : 2
B^10         :
[[  -19    40     0]
 [  -10    21     0]
 [ 1023 -2046  1024]]
column sums  : [1. 1. 1.]
eigen M      : [0.4 0.5 1. ]
x after 10   : [0.5672 0.2661 0.1666]
x after 40   : [0.5667 0.2667 0.1667]
eigen C      : [1. 2.]
eigen C+1e-6 : [0.999 2.001]
cond V(C)    : 2000.0
expm(0.5A)   : 1.9134
formula      : 1.9134
\end{verbatim}
\end{lab}

The first line of that output lists the eigenvalues in increasing order, $1,1,4$. The next three
lines are the residuals $A\mathbf{v}-\lambda\mathbf{v}$ for the three eigenvectors that the library
returned. A residual of order $10^{-16}$ means that the equation is satisfied to the limit of
machine precision; no residual is exactly zero because eigenvectors are normalized and computed
iteratively. The library eigenvectors are also not $(1,1,1)$ and $(1,-1,0)$ but normalized
multiples: the direction is the same, the length differs, and diagonalization cares only about
direction.

The two $A^{10}$ blocks give identical integer entries, with largest difference
$1.05\times10^{-9}$. That difference is not an algebraic error: entries of size
$3.5\times10^{5}$ are computed through $V$, $V^{-1}$, and powers of real numbers, so rounding
errors of order $10^{-16}$ grow along with them.

The value \texttt{cond V} of $1.5227$ is close to one, and that is good news: $A$ is symmetric, so
its eigenvectors can be chosen mutually perpendicular and $P$ is nearly orthogonal.
Diagonalization for such matrices is far more stable.

The power iteration block repeats the computation in Figure~\ref{fig:iterasi-pangkat}: the angle
rises from $0$ to $26.21$ degrees, staying below the limit $26.57$, while the Rayleigh quotient
rises from $4.0000$ to $4.9938$ toward $5$.

The $B$ block shows the other side: its eigenvalues are $1,1,2$ and \texttt{rank(B - I)} equals
$2$, so the nullity is $1$. The geometric multiplicity for $\lambda=1$ is only one, while the
algebraic one is two; the matrix is defective. The $B^{10}$ block shows the consequence: the
entries $-19$ and $1024$ sit side by side, the first growing linearly and the second
exponentially.

The $M$ block verifies that every column sums to one, so that $\mathbf{x}_{k}$ stays stochastic.
Its eigenvalues are $0.4$, $0.5$, and $1$; after forty steps the numbers stop moving at six
digits, matching the rate $(1/2)^{k}$ derived by hand.

The last block gives one warning: the matrix $C$ has eigenvalues $1$ and $2$, and if one zero
entry below the diagonal is replaced by $10^{-6}$, its eigenvalues shift to $0.999$ and $2.001$; a
shift of $10^{-3}$, a thousand times larger than the perturbation. That number can be computed
exactly: the roots of $\lambda^{2}-3\lambda+1.999$ are $1.5\mp\sqrt{2510}/100$, that is
$0.999000998$ and $2.000999002$. The shift is proportional to $10^{-6}\cdot1000/(2-1)$: the
perturbation times the large entry divided by the gap between the two eigenvalues. The condition
number $2000$ captures the same amplification. Eigenvalues are still continuous in the entries,
but they need not be Lipschitz continuous with constant one.

\begin{tantangan}
Find a $2\times2$ matrix whose two eigenvalues are real and distinct, but whose two eigenvectors
are not perpendicular. Show that your choice is correct by computing the dot product, then explain
why a symmetric matrix cannot be used (the answer waits in Chapter 10).

Hint 1: two vectors are perpendicular exactly when their dot product is zero. Start from a simple
upper triangular matrix, then compute both eigenvectors without rounding.

Hint 2: the matrix $\begin{pmatrix}2&1\\0&3\end{pmatrix}$ is enough. The eigenvector for $3$ is
$(1,1)$ and the one for $2$ is $(1,0)$; their dot product is $1$, not zero. Compare with the
symmetric matrix $A$ in the lab, which gives $\texttt{cond V}\approx1.5$.

If you get stuck, go back to the section Similar matrices and invariants, then to the paragraph on
the condition number in Computational lab.
\end{tantangan}

\begin{untukdosen}
A teaching order that has worked: eigenvalues first, as preserved directions, then the
characteristic polynomial as a computational tool. If the polynomial comes first, students see
roots of a polynomial and forget that what is being sought is a direction. After eigen spaces,
introduce geometric multiplicity through nullity, then the condition for diagonalization. The
shear example should come before the condition theorem, so that the theorem answers a question
that is already felt.

Two typical misconceptions. First, students assume that a repeated eigenvalue always means not
diagonalizable. The matrix $(2,1,1;1,2,1;1,1,2)$ in the example section refutes that with
geometric multiplicity two. Second, they write eigenvectors as rows of $P$, then obtain a
$P^{-1}AP$ that is not diagonal and blame the formula. The rule to repeat: eigenvectors are
columns, and the column order of $P$ must be the same as the diagonal order of $D$.

Class activity: split the class into two groups for the matrix $(2,1,1;1,2,1;1,1,2)$. The first
group computes $A^{10}$ through diagonalization, the second through repeated powers or fast
repeated multiplication. Compare the results on the board, then ask the class to explain why the
two roads must meet. The noise usually arises over the denominator $3$ in $P^{-1}$.

Sign of readiness for what follows: students can say, without computing, why the shear matrix is
not diagonalizable, and can guess the form of $S^{k}$ for a $2\times2$ shear. If they mention the
appearance of a factor $k$, they are ready for Chapter 10.
\end{untukdosen}

\section{Summary}
\tingkat{basic}

\begin{ringkasan}
\begin{itemize}
  \item $\lambda$ is an eigenvalue of $A$ if there is $\mathbf{v}\neq\mathbf{0}$ with
        $A\mathbf{v}=\lambda\mathbf{v}$. Eigenvalues are the roots of
        $p_{A}(\lambda)=\det(\lambda I-A)$, and $E_{\lambda}=\operatorname{null}(A-\lambda I)$ is
        the corresponding subspace.
  \item The degree of $p_{A}$ is $n$, so there are at most $n$ distinct eigenvalues. Over
        $\mathbb{C}$ there are always $n$ when counted with multiplicity; over $\mathbb{R}$ there
        need not be, as with the rotation matrix.
  \item $\operatorname{tr}(A)=\sum_{i}\lambda_{i}$ and $\det(A)=\prod_{i}\lambda_{i}$, both
        counted with multiplicity. This is a cheap check for every hand computation.
  \item For every eigenvalue, $1\le\dim E_{\lambda}\le$ the algebraic multiplicity. Eigenvectors
        with distinct eigenvalues are always linearly independent.
  \item $A$ is diagonalizable exactly when there are $n$ independent eigenvectors, and exactly
        when the algebraic multiplicity equals the geometric one for every eigenvalue. Otherwise
        $A$ is called defective, and $A^{k}$ in general contains a factor $k$, not only
        $\lambda^{k}$.
  \item Similar matrices have the same characteristic polynomial, eigenvalues, trace, and
        determinant; the converse does not hold. Eigenvectors change along with the basis.
\end{itemize}
\end{ringkasan}

\section{Exercises}
\tingkat{intermediate}

\begin{latihan}
\begin{enumerate}
  \item Determine the eigenvalues and eigenvectors of
        $A_{1}=\begin{pmatrix}3&1\\ 2&2\end{pmatrix}$, then assemble $P$ and $D$ and check
        $P^{-1}A_{1}P=D$ by multiplying back.

  \item Compute $G^{10}$ for $G=\begin{pmatrix}4&2\\ 1&3\end{pmatrix}$ through diagonalization,
        then check the result with the NumPy function \texttt{matrix\_power}. Write down the
        general form of $G^{k}$ as well.

  \item Prove that if $\lambda$ is an eigenvalue of $A$ then $\lambda^{2}$ is an eigenvalue of
        $A^{2}$, and in general $\lambda^{k}$ is an eigenvalue of $A^{k}$ for every integer
        $k\ge1$. Explain why the converse fails: give an $A$ with $A^{2}=-I$, then compare the
        eigenvalues of $A^{2}$ with the real eigenvalues of $A$.
        (Hint: use the rotation matrix from the start of the chapter; it has no real eigenvalue.)

  \item Prove that $A$ is invertible if and only if $0$ is not an eigenvalue of $A$, and that in
        that case the eigenvalues of $A^{-1}$ are $1/\lambda$ for every eigenvalue $\lambda$ of
        $A$.
        (Hint: start from $A\mathbf{v}=\lambda\mathbf{v}$, multiply on the left by $A^{-1}$, and
        make sure that $\lambda\neq0$ so that the division is legitimate.)

  \item Assess this claim: \emph{two square matrices of the same size that have the same
        characteristic polynomial are necessarily similar.} Decide whether it is true, false, or
        not supported by enough evidence, and include a reason or a counterexample.

  \item Prove that $A$ and $A^{\top}$ have the same characteristic polynomial, using
        $\det(\lambda I-A)=\det((\lambda I-A)^{\top})$ and the properties of the determinant from
        Chapter 8. Explain why this means that the transpose does not change the eigenvalues but
        usually changes the eigenvectors.

  \item The matrix $K=\begin{pmatrix}2&1&0\\ 0&2&0\\ 0&0&1\end{pmatrix}$ has eigenvalues $2,2,1$.
        Determine the rank of $K-2I$, then decide whether $K$ is diagonalizable. Justify your
        answer with the multiplicity condition.

  \item A two-state Markov chain has matrix
        $T=\begin{pmatrix}0.9&0.4\\ 0.1&0.6\end{pmatrix}$. Determine the eigenvalues of $T$, its
        fixed point, and the number of steps needed for the initial error to shrink below
        $10^{-6}$. Check your answer with Python.
        (Hint: once the eigenvalues are known, the error shrinks like $|\lambda_{2}|^{k}$; solve
        $|\lambda_{2}|^{k}<10^{-6}$ with logarithms.)

  \item Prove: if $A$ is diagonalizable with $A=PDP^{-1}$, then
        $\operatorname{tr}(A^{k})=\sum_{i}\lambda_{i}^{k}$ for every $k\ge1$, using
        $\operatorname{tr}(XY)=\operatorname{tr}(YX)$. For the defective $B$ above, compute
        $\operatorname{tr}(B^{10})$ from the printed matrix, then compare it with
        $1^{10}+1^{10}+2^{10}$. Explain why that equality still holds although $B$ is defective,
        and at which point the proof above used diagonalization.
        (Hint: the trace does not change under a change of basis; use the form of $B^{k}$ already
        derived.)
\end{enumerate}
\end{latihan}

\begin{untukpeneliti}
What is really proved in this chapter are algebraic statements over any field and in finite
dimension: the equivalence of eigenvalues with roots of the characteristic polynomial, the
independence of eigenvectors for distinct eigenvalues, the inequality
$1\le\dim E_{\lambda}\le$ the algebraic multiplicity, and the equivalence of diagonalization with
the two conditions used here. The standard references are Axler, \emph{Linear Algebra Done Right}
4th edition, sections 5A and 5D (5.55, the four equivalent conditions, and 5.54); a matrix version
is in Anton, \emph{Elementary Linear Algebra} 11th edition, sections 5.1 and 5.2.

Two things are more often popularized than proved. First, the claim that eigenvalues ``measure
importance'', the basis of the PageRank explanation: the principal eigenvector of a stochastic
matrix gives steady weights only if the chain is irreducible and aperiodic, and that is the content
of the Perron-Frobenius theorem, which is not proved here; a full proof is in Strang,
\emph{Introduction to Linear Algebra} 6th edition, Appendix 8. Second, the claim that the
characteristic polynomial is a practical way to compute eigenvalues: for $n\ge5$ there is no
general root formula, and libraries use the QR algorithm on matrices, not polynomial factoring.
This chapter also stops at the diagonalizable case; Jordan form and the minimal polynomial are not
discussed, and are in Hoffman and Kunze, \emph{Linear Algebra} 2nd edition, sections 6.3 and 7.3.

Open questions: how far eigenvalues shift when the entries are perturbed, whose quantitative answer
requires the Bauer-Fike theorem or pseudospectra; and whether defectiveness itself is numerically
stable. Limits of validity: finite dimension, real or complex field, with no claim about the
stability of algorithms.
\end{untukpeneliti}

\peta{This chapter follows Anton, \emph{Elementary Linear Algebra} 11th edition, Chapter 5:
section 5.1 (Eigenvalues and Eigenvectors), 5.2 (Diagonalization), 5.3 (Complex Vector Spaces),
5.4 (Application: Differential Equations), and 5.5 (Application: Dynamical Systems and Markov
Chains). A more abstract counterpart is in Axler, \emph{Linear Algebra Done Right} 4th edition,
Chapter 5: 5A (eigenvalues and eigen spaces), 5B (existence of eigenvalues over complex vector
spaces), and 5D (diagonalizable operators, including 5.55). For Markov chains and the
Perron-Frobenius theorem, see Strang, \emph{Introduction to Linear Algebra} 6th edition, sections
6.1, 6.2, 6.5, and Appendix 8. Canonical forms for defective matrices are in Hoffman and Kunze,
\emph{Linear Algebra} 2nd edition, Chapter 6 (Elementary Canonical Forms), section 6.3 for the
minimal polynomial, and Chapter 7 (The Rational and Jordan Forms), section 7.3 for the Jordan
form.}
